\documentclass[10pt, xcolor={dvipsnames}]{beamer}

%========================
% CẤU HÌNH HỆ THỐNG
%========================
\usepackage[utf8]{inputenc}
\usepackage[T5]{fontenc}
\usepackage[vietnamese]{babel}
\usepackage{tikz}
\usetikzlibrary{arrows.meta, positioning, calc, shapes.geometric, shadows, patterns}
\usepackage{amsmath, amssymb, amsthm}
\usepackage{tcolorbox}
\tcbuselibrary{skins, shadows, themes}

%========================
% TÔNG MÀU & STYLE
%========================
\usetheme{Madrid}
\definecolor{mainblue}{RGB}{0, 102, 104}
\definecolor{darkblue}{RGB}{0, 32, 96}
\definecolor{lightorange}{RGB}{255, 140, 0}
\definecolor{softblue}{RGB}{235, 245, 255}

\setbeamercolor{palette primary}{bg=mainblue, fg=white}
\setbeamercolor{palette secondary}{bg=lightorange, fg=white}
\setbeamercolor{structure}{fg=mainblue}

% Hộp nội dung đặc sắc
\newtcolorbox{detailbox}[2]{
    enhanced, colback=white, colframe=#1, title=#2,
    fonttitle=\bfseries, arc=3pt, boxrule=1pt, shadow={1mm}{-1mm}{0mm}{black!10}
}

\title{LÝ THUYẾT ĐỒ THỊ}
\subtitle{Toàn tập: Từ Cơ bản đến Nâng cao}
\author{Nguyễn Như Bảo}
\date{Học kỳ II - 2026}

\begin{document}

\frame{\titlepage}
%===========================================================
% PHẦN 2: LÝ THUYẾT (25 SLIDE)
%===========================================================
\section{Phần 2: Lý thuyết cơ sở}

\begin{frame}{2.1. Phân loại: Đơn đồ thị vô hướng}
    \begin{columns}
        \begin{column}{0.6\textwidth}
            \begin{itemize}
                \item Không có cạnh song song.
                \item Không có khuyên (loop).
                \item Quan hệ đối xứng: $(u,v) \equiv (v,u)$.
            \end{itemize}
        \end{column}
        \begin{column}{0.4\textwidth}
            \centering
            \begin{tikzpicture}
                \node[draw,circle](1){1}; \node[draw,circle,right=of 1](2){2};
                \draw (1)--(2);
            \end{tikzpicture}
        \end{column}
    \end{columns}
\end{frame}

\begin{frame}{2.2. Đa đồ thị (Multigraph)}
    Cho phép nhiều cạnh nối giữa hai đỉnh.
    \centering
    \begin{tikzpicture}
\node[draw,circle](1){1}; \node[draw,circle,right=of 1](2){2};
        \draw (1) to [bend left] (2); \draw (1) to [bend right] (2);
    \end{tikzpicture}
\end{frame}

\begin{frame}{2.3. Giả đồ thị (Pseudograph)}
    Cho phép tồn tại khuyên (cạnh nối 1 đỉnh với chính nó).
    \centering
    \begin{tikzpicture}
        \node[draw,circle](1){1}; \draw (1) to [loop above] (1);
    \end{tikzpicture}
\end{frame}

\begin{frame}{2.4. Đồ thị có hướng (Directed Graph)}
    \begin{itemize}
        \item Cạnh có hướng (cung) xác định từ đỉnh đầu đến đỉnh cuối.
        \item $(u,v) \neq (v,u)$.
    \end{itemize}
    \centering
    \begin{tikzpicture}
        \node[draw,circle](1){1}; \node[draw,circle,right=of 1](2){2};
        \draw[->,>=Stealth] (1)--(2);
    \end{tikzpicture}
\end{frame}

\begin{frame}{2.5. Đồ thị có trọng số (Weighted Graph)}
    Mỗi cạnh gắn với một giá trị số $w(e)$.
    \centering
    \begin{tikzpicture}
        \node[draw,circle](1){A}; \node[draw,circle,right=of 1](2){B};
        \draw (1) -- node[above]{5} (2);
    \end{tikzpicture}
\end{frame}

\begin{frame}{2.6. Khái niệm kề và liên thuộc}
    \begin{itemize}
        \item \textbf{Kề (Adjacent):} Hai đỉnh có cạnh nối chung.
        \item \textbf{Liên thuộc (Incident):} Cạnh nối đỉnh $u$ và $v$ thì liên thuộc với $u$ và $v$.
    \end{itemize}
\end{frame}

\begin{frame}{2.7. Bậc của đỉnh (Degree)}
    \begin{detailbox}{mainblue}{Định nghĩa}
        $d(v)$ là số cạnh liên thuộc với đỉnh $v$.
    \end{detailbox}
    \begin{itemize}
        \item \textbf{Đỉnh treo:} Bậc bằng 1.
        \item \textbf{Đỉnh cô lập:} Bậc bằng 0.
    \end{itemize}
\end{frame}

\begin{frame}{2.8. Định lý bắt tay (Handshaking Lemma)}
    \begin{alertblock}{Định lý}
        $\sum_{v \in V} d(v) = 2|E|$
    \end{alertblock}
    \textit{Hệ quả: Số đỉnh bậc lẻ luôn là số chẵn.}
\end{frame}

\begin{frame}{2.9. Bậc trong đồ thị có hướng}
    \begin{itemize}
        \item \textbf{Bán bậc vào ($d^-$):} Số cung trỏ vào đỉnh.
        \item \textbf{Bán bậc ra ($d^+$):} Số cung trỏ ra từ đỉnh.
        \item \textbf{Tổng:} $\sum d^- = \sum d^+ = |E|$.
    \end{itemize}
\end{frame}

\begin{frame}{2.10. Biểu diễn: Ma trận kề (Adjacency Matrix)}
    Sử dụng mảng 2 chiều $A[n][n]$. 
    $A[i][j] = 1$ nếu có cạnh, ngược lại bằng 0.
\end{frame}

\begin{frame}{2.11. Ưu nhược điểm Ma trận kề}
    \begin{itemize}
        \item \textbf{Ưu:} Kiểm tra kề $O(1)$.
        \item \textbf{Nhược:} Tốn bộ nhớ $O(V^2)$. Không tốt cho đồ thị thưa.
    \end{itemize}
\end{frame}

\begin{frame}{2.12. Biểu diễn: Danh sách kề (Adjacency List)}
Mỗi đỉnh quản lý một danh sách các đỉnh kề nó.
    \textbf{Bộ nhớ:} $O(V+E)$. Rất hiệu quả.
\end{frame}

\begin{frame}{2.13. Đường đi (Path)}
    Dãy các đỉnh $v_0, v_1, ..., v_k$ sao cho $(v_i, v_{i+1})$ là cạnh. 
    Độ dài đường đi = $k$ (số cạnh).
\end{frame}

\begin{frame}{2.14. Chu trình (Cycle)}
    Đường đi có đỉnh đầu trùng đỉnh cuối ($v_0 = v_k$).
\end{frame}

\begin{frame}{2.15. Tính liên thông (Connectivity)}
    Đồ thị liên thông nếu giữa 2 đỉnh bất kỳ luôn có đường đi.
\end{frame}

\begin{frame}{2.16. Thành phần liên thông}
    Các cụm đỉnh tối đại mà trong đó mọi cặp đỉnh đều có đường đi tới nhau.
\end{frame}

\begin{frame}{2.17. Đồ thị đầy đủ $K_n$}
    Mọi cặp đỉnh phân biệt đều có cạnh nối. 
    Số cạnh: $|E| = \frac{n(n-1)}{2}$.
\end{frame}

\begin{frame}{2.18. Đồ thị vòng $C_n$}
    Các đỉnh tạo thành một chu trình duy nhất.
\end{frame}

\begin{frame}{2.19. Đồ thị phân đôi (Bipartite Graph)}
    Tập đỉnh $V$ chia thành 2 tập $U, W$ sao cho cạnh chỉ nối giữa $U$ và $W$.
\end{frame}

\begin{frame}{2.20. Đồ thị phân đôi đầy đủ $K_{m,n}$}
    Mọi đỉnh của $U$ đều nối tới mọi đỉnh của $W$.
\end{frame}

\begin{frame}{2.21. Khái niệm Cây (Tree)}
    Đồ thị liên thông và không có chu trình.
\end{frame}

\begin{frame}{2.22. Tính chất của Cây}
    Một cây có $n$ đỉnh thì luôn có đúng $n-1$ cạnh.
\end{frame}

\begin{frame}{2.23. Rừng (Forest)}
    Đồ thị không có chu trình (mỗi thành phần liên thông là một cây).
\end{frame}

\begin{frame}{2.24. Đồ thị phẳng (Planar Graph)}
    Có thể vẽ trên mặt phẳng mà các cạnh không giao nhau.
\end{frame}

\begin{frame}{2.25. Công thức Euler cho đồ thị phẳng}
    \begin{detailbox}{mainblue}{Định lý}
        $V - E + F = 2$ (Với $F$ là số mặt).
    \end{detailbox}
\end{frame}

\end{document}
